\chapter{Formato dei Dati di Input e Output}

\section{Input}

La struttura \texttt{HashRBTree} accetta in input coppie di stringhe 
\texttt{(chiave, valore)} tramite i metodi pubblici \texttt{insert}, 
\texttt{search} e \texttt{remove}. Non è previsto un formato di file esterno: 
le operazioni vengono invocate direttamente nel codice, come mostrato nel 
file \texttt{main.cpp}.

Ogni operazione richiede i seguenti parametri:

\begin{itemize}
    \item \texttt{insert(key, value)} --- due stringhe: la chiave da inserire 
    e il valore associato;
    \item \texttt{search(key, value)} --- la chiave da cercare (il parametro 
    \texttt{value} è ignorato in questa operazione);
    \item \texttt{remove(key, value)} --- la chiave della coppia da rimuovere 
    (il parametro \texttt{value} è ignorato in questa operazione).
\end{itemize}

\section{Output}

L'output è prodotto su \texttt{stdout} tramite \texttt{std::cout}. 
Le operazioni producono i seguenti messaggi:

\subsection{Search}

\begin{itemize}
    \item Se la chiave è presente:
\begin{lstlisting}
Trovato: [chiave] -> valore
\end{lstlisting}
    \item Se la chiave non è presente:
\begin{lstlisting}
Chiave "chiave" non trovata.
\end{lstlisting}
\end{itemize}

\subsection{Remove}

\begin{itemize}
    \item Se la chiave non è presente:
\begin{lstlisting}
Chiave "chiave" non trovata.
\end{lstlisting}
\end{itemize}

\subsection{Print}

Il metodo \texttt{print()} produce una rappresentazione grafica dell'albero 
con indentazione progressiva. Ogni nodo è stampato nel formato:

\begin{lstlisting}
[id | colore] chiave:valore, chiave2:valore2
\end{lstlisting}

dove \texttt{colore} è \texttt{R} per \texttt{RED} e \texttt{B} per 
\texttt{BLACK}. I simboli \texttt{+--}, \texttt{`--} e \texttt{|} 
indicano la struttura gerarchica dell'albero. Un esempio di output:

\begin{lstlisting}
+-- [1938937894 | B] banana:frutto giallo
    +-- [253337143 | R] apple:frutto rosso
    |   +-- [193491643 | B] fig:frutto mediterraneo
    |   `-- [735886645 | B] elderberry:frutto viola
    `-- [1986069970 | B] cherry:frutto rosso piccolo
        `-- [2090176867 | R] date:frutto dolce
\end{lstlisting}

\subsection{Albero Vuoto}

Se \texttt{print()} viene invocato su un albero privo di elementi:

\begin{lstlisting}
Albero vuoto.
\end{lstlisting}